%% ma: L12-B15-0026
%% lop: 12
%% chuong: 5
%% bai: 15
%% dang: 12.15.21
%% loai: TL
%% muc: TH
%% dap_an: "-"
%% nguon: "(NBV)-10 ĐỀ KIỂM TRA HK2-TOÁN 12-DUNG CHUNG 3 SÁCH.pdf (H04, Đề số 4), Phần 4 Câu 2 (tự luận)"
%% kien_thuc: "Bài 7 (tọa độ trung điểm), Bài 15 (phương trình tham số của đường thẳng đi qua hai điểm)"
%% trang_thai: chinh-thuc
%% ghi_chu: "Đáp án gốc khớp: x=1, y=-2+t, z=2t. Đã duyệt 24/10/2026."
\begin{ex}
Trong không gian $Oxyz$, cho ba điểm $A(1;-2;0)$, $B(2;-1;3)$, $C(0;-1;1)$. Viết phương trình đường trung tuyến $AM$ của tam giác $ABC$.
\loigiai{$M$ là trung điểm $BC$ nên $M(1;-1;2)$. $\overrightarrow{AM}=(0;1;2)$.

Đường trung tuyến $AM$ đi qua $A(1;-2;0)$ với vectơ chỉ phương $\overrightarrow{AM}=(0;1;2)$ nên có phương trình $\begin{cases}x=1\\y=-2+t\\z=2t.\end{cases}$}
\end{ex}
